\chapter{Conclusion and Future Work}
\label{chap:conclusion}

\section{Conclusion}

This thesis developed and evaluated a modular Retrieval-Augmented Generation system for structured product-development data. The implementation carries field and relation policy through retrieval, conversational state, prompt construction, generation, and delivery. Its evaluation uses eight attack families and records internal exposure, raw generation where available, control action, delivered output, integrity compliance, and attack-specific authorised task success as separate observations.

The principal architectural finding is that pre-generation role projection performed strongly under the evaluated direct-confidentiality conditions because disallowed fields were removed before the generator received them. The diagnostic sensitivity-evaluation path deliberately retained labelled protected evidence and exposed a much harder control problem. The historical implementation consequently showed substantial confidentiality and integrity failures across direct values, reconstructed rows, retained state, relations, indexed-record existence, and poisoned content. These outcomes cannot be represented by one generic leakage rate.

The final hardened package recorded no unauthorised attack-specific outcome in the finite matrix. This is a substantial descriptive improvement over the historical package, especially where protected information had previously remained model-visible. Prompt and source equivalence is incomplete and several controls changed together, so the package comparison is not a component-level causal estimate. It applies to the tested workbook, models, prompts, scorers, and sample rather than to RAG systems in general.

Matched experiments provide narrower evidence for controls at relation, state, query-handling, and delivery boundaries. Some vulnerable guards-off conditions produced clear changes in the measured outcome; other off/on comparisons remained zero-to-zero and establish only guard execution or context treatment. A supplemental identical-final-raw-answer study isolated a final-answer verifier effect, with measurable over-triggering. A matched synthetic-trigger study showed prompt-injection integrity enforcement at different stages in the two modes.

The prospectively frozen poisoned-row challenge provides the clearest additional integrity result. Enabling the composite prompt-injection guard reduced delivered exact-canary compliance from 300/300 to 0/300. Secure mode acted through pre-generation quarantine, whereas sensitivity-evaluation mode retained 145/150 raw canary responses with the guard enabled and removed the exact canary before delivery. Protected ingredient-plus-percentage leakage was zero in both arms. The authorised positive control was structurally non-evaluable because the linked formulation was retrieved in 0/300 protected records and never entered model context; utility preservation and guard-caused utility loss therefore remain unmeasured. Separate API generations, fixed execution order, and the composite switch bound the attribution to the configured component under the frozen condition.

\subsubsection*{Answers to the research questions}

\paragraph{RQ1: Policy continuity.}
Field-level, role-based sensitivity policy remained most effective when restrictions were enforced structurally before generation and re-evaluated at later security-relevant boundaries. In the historical system, secure-mode projection prevented several direct disclosures, but policy was not preserved automatically across multi-turn reconstruction, access transitions, relation handling, protected-record existence, and other derived representations. Failures were more frequent when protected evidence was deliberately kept model-visible in sensitivity-evaluation mode. The final hardened package recorded no unauthorised attack-specific outcome in the finite evaluation matrix, but this result is limited to the tested conditions and does not establish general policy preservation.

\paragraph{RQ2: Failure boundaries.}
The historical evaluation located different failures at different pipeline boundaries. A01 and A08 demonstrated retrieval-supported delivered disclosure, while A02 demonstrated partial and complete multi-turn reconstruction across retrieval, focus state, memory, prompt construction, and generation. A03 recorded complete post-downgrade delivery together with complete retained-state and downgraded-retrieval exposure; the historical evidence therefore locates the failure at the post-transition retrieval/state boundary but does not isolate either pathway as the cause of the delivered value. A04 exposed protected relation structure; and A05 revealed protected indexed-record existence. A06 produced a small sensitivity-evaluation-mode confidentiality failure without measured canary compliance, whereas A07 produced complete canary compliance without protected-value disclosure, making it an integrity rather than confidentiality failure. Internal exposure was substantially broader than delivered failure, showing that retrieval or context exposure, retained state, raw generation, integrity compliance, and final delivery must be evaluated separately.

\paragraph{RQ3: Control contribution and task preservation.}
Within the matched and supplemental conditions, measurable attack-specific reductions were demonstrated for the relation-access guard in A04, the membership/existence guard in A05, the embedding-probe guard in A08, the supplemental final-answer verifier in A02, and the prompt-injection guard in the supplemental A06 and A07-S studies. Access-change clearing in A03 removed the measured post-transition state exposure but did not demonstrate a delivered-leakage reduction because both matched arms already delivered zero leakage; the earlier A01 and A02 verifier ablations, the earlier A06 ablation, and the A07-N relation-guard ablation were likewise zero-to-zero for their selected delivered outcomes. The corresponding authorised positive-control tasks were preserved in the component studies supporting the A02, A04, A05, A07-S, and A08 outcome reductions. In the supplemental A06 challenge, however, the authorised positive control was structurally non-evaluable because the linked protected formulation was never retrieved, so neither utility preservation nor guard-caused utility loss can be inferred from that experiment.

These results favour structural enforcement before generation whenever policy permits it. A design that intentionally keeps protected evidence model-visible requires additional domain-specific controls and observable enforcement across every subsequent boundary. Delivery filtering can prevent a scored user-visible outcome while leaving upstream exposure intact. Secure structured RAG therefore depends on policy continuity across retrieval, relations, state, prompt construction, generation, and delivery, with evidence interpreted at the stage affected by each intervention.

\section{Limitations}

The evaluation used one workbook, five targets per attack, \texttt{gpt-4o-mini}, \texttt{all-MiniLM-L6-v2}, retrieval depth five, temperature 0.0, templated prompts, and controlled warm-ups. Each 150-conversation unauthorised cell reuses a small target, template, role, and conversation-length matrix with five API repetitions. These are descriptive observations, not 150 independent attack designs. Adaptive strategies, additional domains and models, cross-session state, distributed caches, concurrent users, external tools, and role changes during generation remain outside scope.

The architecture's modularity supports reuse of some control concepts, while transferability remains partial. Ordered sensitivity labels, role-to-label policy, field-sensitivity mappings, explicit relation policy, role projection, access-change state clearing, separation of retrieval, context, raw-answer, and delivery stages, and stage-level telemetry are configuration-driven or architectural elements. Spreadsheet field mappings, logical field names, identifier extraction patterns, configured relation types, exact restricted-value inventories, lexical probe patterns, attack-specific canaries and triggers, target scorers, and structured entity-construction assumptions remain schema- or attack-specific. A new dataset or substantially different policy would require reconfiguration and likely implementation adaptation, followed by renewed validation of every transferred control; the evaluation did not demonstrate plug-and-play portability.

Exact-value and field-aware scorers may miss semantic, approximate, translated, transformed, or derived disclosure. The replay corpus combines historical delivered outputs, preserved matched raw outputs, and synthetic controls under a constructed class balance; its confusion counts describe the frozen corpus rather than deployment prevalence or real-world false-refusal rates. The supplemental A02 challenge isolates final-answer delivery, shares warm-up generation and pre-final state, and reuses the five thesis targets after renderer development on separate targets. It is therefore not a fully independent target-level or end-to-end utility evaluation.

The supplemental A06 gate checked primary canary vulnerability and reduction but omitted a mandatory authorised retrieval opportunity and task-success threshold. Its authorised control was consequently non-evaluable in both pilot and confirmatory stages. The utility portion should have been stopped or redesigned when that structural floor appeared; the completed evidence supports the primary integrity result only.

Historical package provenance lacks a full source manifest, exact system/API payload, and run-time dataset hash. A01 prompt reconstruction and the complete stored A02 user-prompt sequences improve the comparison, but they do not recover package equivalence. Earlier matched studies captured source snapshots, prompts, configurations, datasets, and pair identifiers, while later supplemental studies added versioned scorers, dependency versions, exact messages, and canonical and serialized index hashes. Dirty working trees and several attack-specific gaps remain. Appendix~\ref{app:provenance-detail} records the complete evidence-package inventory, excluded runs, authoritative digest, and HPC archive limitation.

This progression provides a reproducibility lesson from the study: experiments should capture an immutable source archive or clean commit, a portable dependency lock, dataset and index hashes, complete prompts and API payloads, scorer and guard versions, and result checksums at run time. Later studies implement most of these controls, whereas historical gaps cannot be reconstructed retrospectively.

The interface exposed no raw embedding vectors, complete rankings, similarity scores, or citations; the embedding/rank-framed experiment therefore evaluates retrieval-policy enforcement rather than embedding inversion. The protected-record existence probe concerns the indexed corpus rather than model-training membership. Finite zero-leak observations under recorded attacks and scorers cannot establish resistance to adaptive or unseen attacks.

\section{Future Work}

The implemented system and evidence framework provide a baseline for broader structured-RAG security studies. Immediate extensions include additional generation and embedding models, other structured-data domains, longer and less templated conversations, adaptive red teams, multilingual and transformed disclosure, and unseen prompt-injection variants. Further studies should test alternative policy granularity, different structured datasets and domains, and different relation schemas, with renewed validation of transferred controls. Stochastic studies should use independent seeds and uncertainty estimates.

Control evaluation should also become more granular. Separate switches or factorial designs can distinguish pre-generation quarantine from post-generation artifact filtering, while held-out prompt families can test transfer beyond the current targets. Relation and state experiments can extend to deeper and cyclic graphs, cross-session memory, concurrency, caches, and distributed deployment. The final hardened package offers a concrete comparison point for these studies, not a presumed universal solution.

Utility evaluation requires broader end-to-end benign tasks, corrected authorised retrieval opportunities, and explicit acceptance thresholds for false refusal and scorer-negative preservation. Semantic or context-aware verification should be compared with exact matching, particularly for paraphrases, derived values, and user-supplied content. These extensions would test whether the observed security gains transfer while preserving ordinary use beyond the attack-specific positive controls.
